\documentclass[10pt,a4paper]{amsart}
\usepackage[margin=19mm]{geometry}
\usepackage{amsmath,amssymb,mathtools}
\usepackage[scr=rsfs,cal=euler]{mathalfa}
\usepackage{xfrac}
\usepackage[unicode,hidelinks,bookmarks=false]{hyperref}
\newcommand{\ie}{i.e.\ }
\newcommand{\cf}{cf.\ }
\newcommand{\resp}{respectively\ }
\newcommand{\Subsection}{\subsection}
\providecommand{\nobreakdash}{\nobreak\hbox{-}}
\setlength{\emergencystretch}{2em}
\multlinegap0pt
\usepackage{amssymb}
\usepackage{mathtools}
\usepackage[shortlabels]{enumitem}
\newtheorem{lemma}{Lemma}
\newcommand{\F}{\mathbb F_2}
\newcommand{\W}{\mathcal W}
\DeclareMathOperator{\con}{con}
\newcommand{\eqid}{\approx}
\newcommand{\ivec}[1]{\mathbf e(#1)}
\title{The variety generated by all semirings of order three is nonfinitely based}
\author{AIFA WANG, WENHAO JU, LILI WANG}
\date{Source: arXiv:2609.15703v1 (CC BY 4.0)}
\begin{document}
\maketitle

\noindent\emph{Source and licence.} The variety generated by all semirings of order three is nonfinitely based. Version arXiv:2609.15703v1 (CC BY 4.0), distributed by arXiv under the Creative Commons Attribution 4.0 International licence, http://creativecommons.org/licenses/by/4.0/, as recorded in the arXiv licence metadata for this exact version and shown on the abstract page https://arxiv.org/abs/2609.15703v1. Full licence notice: \texttt{licenses/arxiv-2609.15703-CC-BY-4.0.txt}.\par\smallskip

\noindent\textit{Editorial scope.} Counted unit: the article's lemma lem:isolation (source0.tex lines 797--806). The article's complete original proof of it is delivered with it (source0.tex lines 808--905, one \texttt{proof} environment of 98 lines). No mathematics is written, repaired or re-derived: the statement and the proof are the article's own lines, in the article's own notation, and no retained line is substituted or re-flowed.  Retained uncounted as the source-local context the proof needs: lines 400--410; lines 511--513.  Nothing was omitted from any retained span, so the package carries no omission record.  No retained line contains a citation, so the wrapper carries no bibliography and no bibliography entry.  No retained line contains a figure environment, an image-inclusion command or drawing code, so no image is delivered and the package declares no asset dependency.  Editorial formatting only, in addition: an AMS \texttt{amsart} preamble that restates the article's own declarations and macros used by the retained lines, namely the environments \texttt{lemma} on separate counters, together with the standard packages those lines need.  The retained environments print in this document as Lemma 1 in order of appearance, whereas the article prints the same items with the numbering of its own full document.  The complete native source of the article as distributed in the arXiv e-print for this exact version is bundled unchanged as \texttt{sources/210442/source0.tex}.
\begin{lemma}\label{lem:tests}
Let $u$ be a linear polynomial and let $r$ be a commutative word.
If $\W\models r\preceq u$, then the following conditions hold:
\begin{enumerate}
\item $r$ is linear and $\con(r)\subseteq\con(u)$;
\item every two distinct variables of $r$ occur together in a word of $u$;
\item some word of $u$ has content contained in $\con(r)$;
\item $\ivec{r}$ belongs to the affine span over $\F$ of the incidence
vectors of the words of $u$.
\end{enumerate}
\end{lemma}

\begin{lemma}\label{lem:test-valid}
The variety $\W$ satisfies $\sigma_n$ for every odd $n\geq3$.
\end{lemma}

\begin{lemma}[Isolation]\label{lem:isolation}
Let $n\geq5$ be odd. If $\W\models r\preceq U_n$ for a
commutative word $r$, then
\[
r\in\{p_n,a_n,b_1,\ldots,b_n,q_n\}.
\]
Moreover, the equivalence class of $U_n$ under the identities of
$\W$, in the free commutative ai-semiring, is exactly
$\{U_n,U_n+q_n\}$.
\end{lemma}

\begin{proof}
We first determine the possible words below $U_n$ and then show
that every word originally present in $U_n$ is indispensable.

By Lemma~\ref{lem:tests}, $r$ is linear and uses only variables of
$U_n$. Form a graph $\Gamma_n$ with these variables as vertices,
joining two distinct vertices when they occur together in a word
of $U_n$. Conditions (2) and (3) of that lemma say that
$\con(r)$ is a clique containing the content of at least one word
of $U_n$.

The $x$-vertices form a clique, and $z$ is adjacent to every
$x$- and $y$-vertex. The only $y$-neighbours of $x_i$ are
$y_i,y_{i+1}$. Among the $x$-vertices, the neighbours of $y_i$
are precisely $x_{i-1},x_i$; among the $y$-vertices, its neighbours
are $y_{i-1},y_{i+1}$. We use this description in the following
three cases.

\emph{Case 1: $\con(p_n)\subseteq\con(r)$.}
No $y_i$ is adjacent to every $x$-vertex, since $n\geq5$ and
$y_i$ has only two $x$-neighbours. Hence the only additional
vertex that can occur in the clique $\con(r)$ is $z$.
Linearity now gives $r=p_n$ or $r=q_n$.

\emph{Case 2: $\con(b_i)\subseteq\con(r)$ for some $i$.}
The common $x$-neighbour of $y_i,y_{i+1}$ is only $x_i$.
Also, no additional $y$-vertex is adjacent to both $y_i$ and
$y_{i+1}$ in a cycle of length at least five. Thus the clique
$\{x_i,y_i,y_{i+1},z\}$ is maximal, and $r=b_i$.

\emph{Case 3: $\{x_1,z\}\subseteq\con(r)$.}
Any $y$-variables in $r$ must lie in $\{y_1,y_2\}$, because
they must be adjacent to $x_1$. Every incidence vector of a
word of $U_n$ has an even sum of its $y$-coordinates.
This property is preserved by linear combinations over $\F$,
so Lemma~\ref{lem:tests}(4) gives the same property for
$\ivec{r}$. Thus either both $y_1,y_2$ occur or neither occurs.
If both occur, their only common $x$-neighbour is $x_1$,
and no further $y$-variable is allowed. Hence $r=b_1$.

Suppose that no $y$-variable occurs in $r$. By affine-span
membership, there are $\alpha,\beta,\gamma_1,\ldots,\gamma_n\in\F$
such that
\begin{equation}\label{eq:affine}
\begin{aligned}
\ivec{r}
 &=\alpha\ivec{p_n}+\beta\ivec{a_n}
       +\sum_{i=1}^{n}\gamma_i\ivec{b_i},\\
\alpha+\beta+\sum_{i=1}^{n}\gamma_i&=1.
\end{aligned}
\end{equation}
The $y_i$-coordinate on the left is zero. On the right, it is
$\gamma_{i-1}+\gamma_i$, with $\gamma_0=\gamma_n$.
Thus
\[
\gamma_1=\gamma_2=\cdots=\gamma_n=\gamma.
\]
Since $n$ is odd, $\sum_i\gamma_i=\gamma$ in $\F$.
Since $z$ occurs in $r$, comparison of the $z$-coordinates
and the coefficient condition in \eqref{eq:affine} give
\[
\beta+\gamma=1,\qquad \alpha+\beta+\gamma=1.
\]
It follows that $\alpha=0$. If $(\beta,\gamma)=(1,0)$,
then $\ivec{r}=\ivec{a_n}$ and $r=a_n$.
If $(\beta,\gamma)=(0,1)$, then
$\ivec{r}=\sum_i\ivec{b_i}$.
In this sum every $y_i$ occurs twice, every $x_i$ once, and
$z$ occurs an odd number of times. Hence this vector is
$\ivec{q_n}$ and $r=q_n$. This completes the list of possible
words, because condition (3) of Lemma~\ref{lem:tests} ensures
that at least one of the three cases applies.

Now let $\W\models v\eqid U_n$. Each word of $v$ is below
$v$ and therefore below $U_n$, so every word of $v$ belongs
to the displayed list. We claim that $p_n,a_n,b_1,\ldots,b_n$
must all occur in $v$.

If $p_n$ were absent, assign $1$ to all $x_i$ and $0$ to
$z$ and all $y_i$ in $D$. Then $U_n$ has value $1$,
whereas every remaining listed word contains $z$ and has value
$0$. If $a_n$ were absent, assign $1$ only to $x_1,z$ in
$D$. The word $a_n$ has value $1$, but every other listed
word contains a variable assigned $0$. Either omission would
contradict $v\eqid U_n$.

Finally, if $b_i$ were absent, assign $a$ to
$y_i,y_{i+1}$ and $1$ to all other variables in $H$.
The word $b_i$ has value $\infty$, while every other listed
word contains at most one of $y_i,y_{i+1}$ and has value
$1$ or $a$. Thus $U_n$ and $v$ would again have different
values. This proves the claim.

Consequently, $v$ is either $U_n$ or $U_n+q_n$.
Both possibilities do occur: the first trivially, and the
second by Lemma~\ref{lem:test-valid}. Therefore the equivalence class has
exactly the two asserted members.
\end{proof}

\end{document}
